% Options for packages loaded elsewhere
\PassOptionsToPackage{unicode}{hyperref}
\PassOptionsToPackage{hyphens}{url}
\PassOptionsToPackage{space}{xeCJK}
\documentclass[
  chinese,
]{article}
\usepackage{xcolor}
\usepackage[margin=2cm]{geometry}
\usepackage{amsmath,amssymb}
\setcounter{secnumdepth}{-\maxdimen} % remove section numbering
\usepackage{iftex}
\ifPDFTeX
  \usepackage[T1]{fontenc}
  \usepackage[utf8]{inputenc}
  \usepackage{textcomp} % provide euro and other symbols
\else % if luatex or xetex
  \usepackage{unicode-math} % this also loads fontspec
  \defaultfontfeatures{Scale=MatchLowercase}
  \defaultfontfeatures[\rmfamily]{Ligatures=TeX,Scale=1}
\fi
\usepackage{lmodern}
\ifPDFTeX\else
  % xetex/luatex font selection
  \ifXeTeX
    \usepackage{xeCJK}
    \setCJKmainfont[]{Microsoft YaHei}
  \fi
  \ifLuaTeX
    \usepackage[]{luatexja-fontspec}
    \setmainjfont[]{Microsoft YaHei}
  \fi
\fi
% Use upquote if available, for straight quotes in verbatim environments
\IfFileExists{upquote.sty}{\usepackage{upquote}}{}
\IfFileExists{microtype.sty}{% use microtype if available
  \usepackage[]{microtype}
  \UseMicrotypeSet[protrusion]{basicmath} % disable protrusion for tt fonts
}{}
\makeatletter
\@ifundefined{KOMAClassName}{% if non-KOMA class
  \IfFileExists{parskip.sty}{%
    \usepackage{parskip}
  }{% else
    \setlength{\parindent}{0pt}
    \setlength{\parskip}{6pt plus 2pt minus 1pt}}
}{% if KOMA class
  \KOMAoptions{parskip=half}}
\makeatother
\usepackage{longtable,booktabs,array}
\usepackage{calc} % for calculating minipage widths
% Correct order of tables after \paragraph or \subparagraph
\usepackage{etoolbox}
\makeatletter
\patchcmd\longtable{\par}{\if@noskipsec\mbox{}\fi\par}{}{}
\makeatother
% Allow footnotes in longtable head/foot
\IfFileExists{footnotehyper.sty}{\usepackage{footnotehyper}}{\usepackage{footnote}}
\makesavenoteenv{longtable}
\usepackage{graphicx}
\makeatletter
\newsavebox\pandoc@box
\newcommand*\pandocbounded[1]{% scales image to fit in text height/width
  \sbox\pandoc@box{#1}%
  \Gscale@div\@tempa{\textheight}{\dimexpr\ht\pandoc@box+\dp\pandoc@box\relax}%
  \Gscale@div\@tempb{\linewidth}{\wd\pandoc@box}%
  \ifdim\@tempb\p@<\@tempa\p@\let\@tempa\@tempb\fi% select the smaller of both
  \ifdim\@tempa\p@<\p@\scalebox{\@tempa}{\usebox\pandoc@box}%
  \else\usebox{\pandoc@box}%
  \fi%
}
% Set default figure placement to htbp
\def\fps@figure{htbp}
\makeatother
\ifLuaTeX
\usepackage[bidi=basic,shorthands=off,,provide=*]{babel}
\else
\usepackage[bidi=default,shorthands=off,,provide=*]{babel}
\fi
\ifLuaTeX
  \usepackage{selnolig} % disable illegal ligatures
\fi
\setlength{\emergencystretch}{3em} % prevent overfull lines
\providecommand{\tightlist}{%
  \setlength{\itemsep}{0pt}\setlength{\parskip}{0pt}}
\usepackage{bookmark}
\IfFileExists{xurl.sty}{\usepackage{xurl}}{} % add URL line breaks if available
\urlstyle{same}
\hypersetup{
  pdflang={zh-CN},
  hidelinks,
  pdfcreator={LaTeX via pandoc}}

\author{}
\date{}

\begin{document}

\section{Embodied RL Roadmap：从经典控制到 Diffusion/VLA
后训练}\label{embodied-rl-roadmapux4eceux7ecfux5178ux63a7ux5236ux5230-diffusionvla-ux540eux8badux7ec3}

\textbf{资料截止：2026-09-12}

\subsection{摘要}\label{ux6458ux8981}

具身智能中的强化学习（reinforcement learning,
RL）经历了从状态/动作空间中的经典控制，到深度策略、仿真到现实（sim-to-real）、真实机器人在线/离线
RL，再到扩散策略、视觉语言动作模型（VLA）后训练的演化。主线并不是``PPO
换成了某个新算法''，而是三件事一起变化：反馈从手工奖励扩展为示范、偏好、任务成功与世界模型反馈；策略从单峰高斯扩展到动作序列的多模态生成；数据从少量人工轨迹扩展为仿真、离线数据和真实机器人数据飞轮。

本文用 Flow-GRPO 作为扩散模型在线 RL 的前置，用 DiffusionNFT 作为
forward-process
后训练的代表，并分析这些图像生成方法迁移到具身动作生成时哪些结构可以复用、哪些地方必须重新设计。图像生成的奖励提升不能直接推出机器人闭环控制提升，这一限制贯穿全文。

\subsection{0.
给初学者的最短预备课}\label{0-ux7ed9ux521dux5b66ux8005ux7684ux6700ux77edux9884ux5907ux8bfe}

可以先把 RL
想成``让一个会试错的控制器学会做事''。机器人看到相机画面和自身关节状态，这些信息组成状态
\(s_t\)；控制器输出动作
\(a_t\)，例如关节速度或夹爪开合；环境执行动作后给出下一个状态
\(s_{t+1}\) 和奖励
\(r_t\)。一段从开始到结束的交互记录叫作轨迹（trajectory）。如果机器人把杯子放到目标位置，最后可能得到较高奖励；如果撞倒杯子，则得到低奖励或惩罚。

这里最容易混淆的三个词是：**策略（policy）**是``看到状态后如何选动作''的规则，**奖励（reward）**是某一步得到的即时分数，**回报（return）**是从当前时刻往后的折扣奖励总和。RL
的学习目标不是记住某一个动作，而是让策略在未来许多步之后仍能取得较高回报。因此，长任务中的``信用分配''很重要：如果
20 步之后才知道任务失败，算法要判断前面哪几步导致了失败。

具身智能又比游戏或图像生成多一层困难：动作真的会改变世界。图像模型可以一次生成很多候选图片，再由评分器打分；机器人每次探索都要花时间，还可能碰撞、磨损或需要人工重置。后文出现``在线''时，指策略与环境持续交互；出现``离线''时，指只从已经收集好的轨迹数据学习。

\subsection{1.
先建立统一坐标系}\label{1-ux5148ux5efaux7acbux7edfux4e00ux5750ux6807ux7cfb}

在普通 MDP 中，策略 \(\pi(a_t|s_t)\) 根据状态选择动作，奖励 \(r_t\)
评价短期行为，价值函数负责长时序信用分配。具身任务额外包含视觉观测、动力学、接触、延迟、安全约束和硬件磨损。因此同一个``在线
RL''在图像生成和机器人控制中并不等价：图像可以在离线评分器上快速反复采样，机器人每一次探索都可能消耗时间、撞击物体或损坏硬件。

本文用六个维度比较方法：反馈来源、策略表示、训练环境、在线程度、信用分配粒度和部署约束。

\subsection{2. 发展脉络}\label{2-ux53d1ux5c55ux8109ux7edc}

\subsubsection{2.1 经典 RL
与机器人控制：从值函数到策略优化}\label{21-ux7ecfux5178-rl-ux4e0eux673aux5668ux4ebaux63a7ux5236ux4eceux503cux51fdux6570ux5230ux7b56ux7565ux4f18ux5316}

早期机器人 RL 主要处理低维状态、动作和可设计奖励。Q-learning
学习``在某个状态做某个动作值不值得''，actor-critic
同时学习策略和价值估计，policy gradient 直接调整策略参数，PPO
则用裁剪后的更新限制策略一次变化太大。它们解决的是不同的优化稳定性问题，但都需要回答同一个问题：怎样从奖励反推出哪些动作值得重复。

真正进入复杂机器人后，难点从``有没有 RL
算法''变成样本效率、探索安全和奖励塑形。比如``把物体放进盒子''可以把成功作为稀疏奖励，但机器人可能很久都碰不到盒子；把距离、接触力和姿态加入奖励会更容易学，却可能诱导策略钻奖励函数的空子。这就是为什么具身
RL 通常把示范学习、课程学习、奖励模型和安全约束放在同一个系统里。

\subsubsection{2.2 深度 RL 与
Sim-to-Real：把探索移入仿真}\label{22-ux6df1ux5ea6-rl-ux4e0e-sim-to-realux628aux63a2ux7d22ux79fbux5165ux4effux771f}

深度视觉策略、域随机化、系统辨识和残差控制把大部分探索移到仿真中。其典型流程是：在仿真中用大量并行环境训练策略，再通过视觉、动力学和执行器随机化迁移到真实机器人。域随机化不是把仿真做得``看起来漂亮''，而是主动改变光照、相机噪声、摩擦、质量、延迟和控制频率，让策略学到对这些变化不敏感的规律。

这条路线的关键假设是：真实世界落在训练时覆盖的变化范围内。优点是样本量可扩展，缺点是仿真偏差会集中出现在接触、摩擦、遮挡和长时序任务上。一个在仿真中成功率很高的策略，仍可能因为真实夹爪略有偏差而抓不稳。因此
sim-to-real 往往还需要少量真实数据做校准或残差修正。

\subsubsection{2.3 真实世界
RL：反馈闭环成为系统工程}\label{23-ux771fux5b9eux4e16ux754c-rlux53cdux9988ux95edux73afux6210ux4e3aux7cfbux7edfux5de5ux7a0b}

The Ingredients of Real-World Robotic Reinforcement Learning
将问题拆成奖励、探索、数据采集、重置、硬件安全和持续训练等组件。这里的``ingredients''可以理解为一套运行条件：奖励要能自动计算，机器人要能在失败后回到可工作的初始状态，数据要记录传感器和动作时间戳，控制器要有速度/力矩边界，训练过程还要能筛掉异常轨迹。

真实世界 RL
由此从单一算法问题转为完整系统问题。自动重置、并行机器人、失败恢复和数据过滤与优化器同样重要。很多论文的性能提升来自整个数据闭环，而不是某一个损失函数单独带来的提升；阅读结果时要把``算法贡献''和``系统工程贡献''分开。

\pandocbounded{\includegraphics[width=0.78\linewidth,height=0.32\textheight,keepaspectratio,alt={真实世界机器人 RL 的原论文素材，来自 Ingredients of Real-World Robotic RL 的源码图}]{figures/ingredients-real-world-rl_asset_1.png}}

\subsubsection{2.4 离线 RL
与示范学习：先从已有经验中学习}\label{24-ux79bbux7ebf-rl-ux4e0eux793aux8303ux5b66ux4e60ux5148ux4eceux5df2ux6709ux7ecfux9a8cux4e2dux5b66ux4e60}

离线 RL
通过固定数据集学习策略，避免一开始就在真实世界探索。数据可以来自专家示范、旧策略、遥操作或过去的失败尝试。D4RL、IQL、CQL、AWAC
和 diffusion-based offline RL
代表了不同路径。离线数据让策略学习更安全，但也带来两个初学者需要特别注意的问题：第一，数据里没有出现过的动作很难可靠评估；第二，数据只记录了发生过的结果，不能自动告诉模型``如果当时换个动作会怎样''。

机器人实践通常采用``示范/离线预训练 +
少量在线改进''，而不是纯离线或从零在线。前者提供一个不会立刻摔倒的初始策略，后者利用真实奖励修正示范中的盲点。

\subsubsection{2.5 Diffusion
Policy：把动作看成多模态序列}\label{25-diffusion-policyux628aux52a8ux4f5cux770bux6210ux591aux6a21ux6001ux5e8fux5217}

Diffusion Policy
将动作序列作为条件生成对象，使用视觉观测和历史动作预测一段 action
chunk。这里的``扩散''可以先用一个直观比喻理解：训练时逐渐给真实动作加噪声，模型学习如何一步步去噪；执行时从噪声开始，反过来逐步生成一段动作。因为生成的是分布而不是一个平均动作，它可以保留多个同样合理的行为模式。

例如，同一个物体可以从左边抓，也可以从右边抓。普通回归模型容易把两条轨迹平均成一条不可执行的中间轨迹；扩散策略可以把它们作为两个模式保留下来。代价是执行时要做多步采样，而且每次生成的动作仍需放回真实闭环中检查。

\pandocbounded{\includegraphics[width=0.78\linewidth,height=0.32\textheight,keepaspectratio,alt={Diffusion Policy 原论文 teaser，来自论文源码}]{figures/rendered/diffusion_policy_teaser.png}}

Diffusion Policy 首先是策略表示和监督/模仿学习方法，不自动等于
RL。要把它用于
RL，需要定义任务奖励、探索机制、动作序列信用分配和安全约束。Diffusion
Policies as an Expressive Policy Class for Offline Reinforcement
Learning 进一步说明扩散分布可以作为离线 RL
的策略类，但仍受数据覆盖和价值外推影响。

\subsection{3. Flow-GRPO：DiffusionNFT
的必要前置}\label{3-flow-grpodiffusionnft-ux7684ux5fc5ux8981ux524dux7f6e}

\subsubsection{3.1 GRPO 的直觉}\label{31-grpo-ux7684ux76f4ux89c9}

对同一个条件 \(c\) 采样 \(N\) 个结果 \(x^{(i)}\)，计算每个结果的奖励
\(R^{(i)}\)，再用组内均值和标准差得到相对优势：

\[A^{(i)}=\frac{R^{(i)}-\operatorname{mean}(R)}{\operatorname{std}(R)+\epsilon}.\]

高于组平均的结果被加强，低于组平均的结果被削弱。它不需要单独训练一个
value model，因此适合把奖励直接用于生成模型后训练。

举一个很小的例子：同一个指令生成 4 张图，评分为 0.2、0.4、0.7、0.8。GRPO
不需要知道``绝对的 0.8
是否足够好''，只需要知道后两张比组内平均更好，于是提高产生它们的概率。对机器人而言，同一观测也可以采样多段候选动作，用任务成功、距离或碰撞惩罚进行组内比较。难点在于机器人奖励通常延迟到整段动作结束，信用分配会比单张图像更难。

\subsubsection{3.2 为什么 flow matching
需要额外处理}\label{32-ux4e3aux4ec0ux4e48-flow-matching-ux9700ux8981ux989dux5916ux5904ux7406}

许多 flow-matching 采样器由确定性 ODE 描述。确定性路径没有传统 policy
gradient 所需的随机探索，也不直接给出每一步的动作概率。Flow-GRPO
的关键步骤是把 ODE 转换为保持边际分布的
SDE，使每一步拥有可计算的随机转移；随后对同一 prompt
的多条生成轨迹计算组相对优势，并把最终图像奖励回传到去噪过程。

\pandocbounded{\includegraphics[width=0.78\linewidth,height=0.32\textheight,keepaspectratio,alt={Flow-GRPO 原论文方法图，来自论文源码渲染}]{figures/rendered/flow_grpo_method.png}}

\pandocbounded{\includegraphics[width=0.78\linewidth,height=0.32\textheight,keepaspectratio,alt={Flow-GRPO 原论文 teaser，来自论文源码渲染}]{figures/rendered/flow_grpo_teaser.png}}

Flow-GRPO 的优点是首次把在线策略梯度式 RL 连接到 flow-matching
生成模型；代价是仍要处理反向采样轨迹、概率比、solver
约束和逐步信用分配。论文报告其在文本到图像任务上改善组合关系、文字渲染和偏好对齐，但这些指标属于图像生成证据，不是具身控制证据。

可以把它理解成两层循环：外层对同一个条件采样多条完整生成路径并比较奖励，内层沿每条路径的多个去噪步骤更新模型。外层回答``哪张结果更好''，内层回答``这条结果中的哪些随机转移应该更常发生''。Flow-GRPO
的工程困难正来自这条内层路径：需要保留或重建每一步的随机转移，并让训练时的概率计算和实际采样
solver 对得上。

\subsection{4. DiffusionNFT：从 reverse trajectory 到 forward
process}\label{4-diffusionnftux4ece-reverse-trajectory-ux5230-forward-process}

\subsubsection{4.1 方法定位}\label{41-ux65b9ux6cd5ux5b9aux4f4d}

DiffusionNFT（Diffusion Negative-aware Fine-Tuning）提出在 forward
diffusion / flow-matching 过程中直接做在线
RL。它把生成结果按奖励划分为正样本和负样本，用二者的预测误差构造隐式的策略改进方向，并把
RL 信号整合进类似监督学习的 velocity objective。

\pandocbounded{\includegraphics[width=0.78\linewidth,height=0.32\textheight,keepaspectratio,alt={DiffusionNFT 方法图，来自论文源码渲染}]{figures/rendered/diffusionnft_method.png}}

与 Flow-GRPO 的对照可以概括为：

\begin{longtable}[]{@{}lll@{}}
\toprule\noalign{}
维度 & Flow-GRPO & DiffusionNFT \\
\midrule\noalign{}
\endhead
\bottomrule\noalign{}
\endlastfoot
优化位置 & reverse sampling trajectory & forward/noising 与
flow-matching objective \\
探索形式 & ODE-to-SDE 后的随机轨迹 & 生成样本的正负对比 \\
轨迹存储 & 需要逐步采样信息 & 只需 clean samples 参与训练 \\
solver & 对采样器有更强约束 & 支持黑盒 solver \\
CFG & 需要处理 CFG/参考策略 & 论文强调 CFG-free 设定 \\
直观理解 & 对轨迹做 policy gradient & 用奖励排序改变 velocity field \\
\end{longtable}

\pandocbounded{\includegraphics[width=0.78\linewidth,height=0.32\textheight,keepaspectratio,alt={DiffusionNFT 性能对比图，来自论文源码渲染}]{figures/rendered/diffusionnft_performance.png}}

论文报告在 SD3.5-Medium 上相对 Flow-GRPO 的效率优势，例如 GenEval 从
0.24 提升到 0.98 的对比；这些数字应理解为图像生成实验中的 paper
fact，不能直接外推到机器人成功率。

对初学者来说，DiffusionNFT 的核心变化可以浓缩成一句话：Flow-GRPO
把``整条反向生成轨迹''当作 RL 的对象，DiffusionNFT
则尽量把奖励信息转换成 forward-process
上的训练信号。模型仍然需要生成样本并获得奖励，但优化阶段主要比较好样本和坏样本在
velocity prediction
上的差异。这样做减少了对完整采样轨迹和显式似然估计的依赖。

这并不意味着它不需要探索或奖励。它仍然需要候选样本、奖励模型和正负划分；改变的是``如何把奖励变成参数更新''。因此论文中``只需
clean
samples''应理解为训练更新不必保存整条反向轨迹，而不是完全不需要生成过程。

\subsubsection{4.2
迁移到具身动作生成：哪些东西可以复用}\label{42-ux8fc1ux79fbux5230ux5177ux8eabux52a8ux4f5cux751fux6210ux54eaux4e9bux4e1cux897fux53efux4ee5ux590dux7528}

\begin{longtable}[]{@{}ll@{}}
\toprule\noalign{}
图像生成 & 具身智能中的对应物 \\
\midrule\noalign{}
\endhead
\bottomrule\noalign{}
\endlastfoot
prompt & 语言指令、视觉状态和任务上下文 \\
clean image & action chunk 或动作轨迹 \\
image reward & 任务成功率、距离、碰撞、能耗、平滑度 \\
正/负图像 & 成功/失败 rollout 或候选动作 \\
forward noising & 动作扰动、latent/action flow \\
单次生成质量 & 长时序闭环控制表现 \\
\end{longtable}

\textbf{可迁移的结构}包括正负样本比较、forward-process 优化和 solver
解耦。\textbf{必须重做的部分}包括时间信用分配、动作连续性、动力学可行性、在线数据飞轮、碰撞和安全过滤、sim-to-real
校准。机器人动作是带动力学的时间序列，单个最终动作 chunk
的奖励无法自动告诉模型哪一个去噪时刻或哪一步接触导致了失败。

因此，DiffusionNFT
更适合作为``扩散动作策略后训练的候选机制''，而不是已经完成验证的机器人算法。具身实验需要至少增加：仿真
benchmark、真实机器人闭环成功率、不同 embodiment
的迁移、失败率/碰撞率、动作平滑度和数据效率。

一个可执行的迁移实验是：先固定一个已经用示范训练好的 diffusion action
policy，在仿真中为同一观测生成一组 action
chunks；用任务成功、末端距离和碰撞惩罚排序；再比较普通行为克隆、轨迹级
Flow-GRPO 风格更新和 forward-process/NFT
风格更新。只有当三种方法使用相同的采样预算、奖励函数和安全过滤，并在未见过的初始位姿上测试，才能判断
NFT 风格更新是否真的带来数据效率提升。

\subsection{5. 从动作扩散到 VLA
与数据飞轮}\label{5-ux4eceux52a8ux4f5cux6269ux6563ux5230-vla-ux4e0eux6570ux636eux98deux8f6e}

OpenVLA 等开放 VLA 模型把视觉、语言和动作统一到一个策略中。扩散/flow
action head 负责连续动作生成，语言视觉编码器负责任务条件。此时 RL
的角色从``从零学策略''逐渐转为：在示范预训练策略上做任务后训练、在线纠错和奖励对齐。

DexFlyWheel
将灵巧操作的数据采集、失败分析和策略改进组织成可循环的数据框架；RL-100
则将真实机器人数据飞轮、离线 RL
和在线提升组合起来，代表``先利用人类先验，再用真实反馈超越示范''的路线。

\pandocbounded{\includegraphics[width=0.78\linewidth,height=0.32\textheight,keepaspectratio,alt={DexFlyWheel 原论文方法图，来自论文源码渲染}]{figures/rendered/dexflywheel_method.png}}

\pandocbounded{\includegraphics[width=0.78\linewidth,height=0.32\textheight,keepaspectratio,alt={X-DiffVLA 原论文框架图，来自论文源码渲染}]{figures/rendered/x_diffvla_framework.png}}

近期工作开始把 flow/diffusion action head 与跨机器人 embodiment
迁移结合起来，但仍要区分``动作头的表达能力''\,``VLA
的语义泛化''和``真实机器人 RL 的闭环适应''。

\subsection{6. 按阶段比较}\label{6-ux6309ux9636ux6bb5ux6bd4ux8f83}

\begin{longtable}[]{@{}lllll@{}}
\toprule\noalign{}
阶段 & 主要数据 & 策略表示 & RL 的主要作用 & 主要瓶颈 \\
\midrule\noalign{}
\endhead
\bottomrule\noalign{}
\endlastfoot
经典控制 & 状态、少量轨迹 & 值函数/高斯策略 & 学控制规律 &
状态维度和奖励设计 \\
深度 RL + 仿真 & 大规模仿真 rollout & 深度策略 & 探索与控制 &
sim-to-real \\
真实世界 RL & 机器人在线数据 & 深度策略/残差策略 & 适应真实反馈 &
样本效率、安全 \\
离线 RL & 示范与历史数据 & Q/actor/diffusion & 从固定数据改进 & OOD
动作、覆盖不足 \\
Diffusion Policy & 示范动作序列 & 扩散动作分布 & 表达多模态动作 &
推理成本、闭环校正 \\
Flow-GRPO/DiffusionNFT & 生成结果与奖励 & flow/diffusion 生成器 &
奖励对齐与后训练 & 奖励黑客、信用分配 \\
VLA + RL & 视觉语言动作数据飞轮 & foundation policy + action head &
任务后训练、迁移 & 长时序、硬件和泛化 \\
\end{longtable}

\subsection{7.
结论与研究问题}\label{7-ux7ed3ux8bbaux4e0eux7814ux7a76ux95eeux9898}

\begin{enumerate}
\def\labelenumi{\arabic{enumi}.}
\tightlist
\item
  \textbf{反馈闭环比算法名称更关键}：没有可靠的奖励、重置、过滤和失败记录，换成更复杂的
  RL 目标也无法稳定提升机器人能力。
\item
  \textbf{扩散策略解决的是分布表达问题}：它可以表达多模态动作，但在线 RL
  还需要解决动作序列信用分配和安全探索。
\item
  \textbf{Flow-GRPO 与 DiffusionNFT
  是两种互补路线}：前者把轨迹级策略梯度接到 flow
  sampler，后者把奖励排序转成 forward-process 的监督式改进。
\item
  \textbf{图像生成结果不能直接证明具身迁移}：迁移假设必须由仿真、真实机器人和跨
  embodiment 实验验证。
\item
  \textbf{RL-100
  类系统提示了落地路线}：人类示范/预训练策略提供先验，真实机器人数据飞轮提供持续反馈，离线
  RL 提高安全性，在线 RL 完成最后的任务适应。
\end{enumerate}

\subsection{8.
阅读与复现顺序}\label{8-ux9605ux8bfbux4e0eux590dux73b0ux987aux5e8f}

\begin{enumerate}
\def\labelenumi{\arabic{enumi}.}
\tightlist
\item
  Diffusion Policy：理解动作扩散策略表示；
\item
  Flow-GRPO：理解 ODE-to-SDE、组相对优势和轨迹级 RL；
\item
  DiffusionNFT：理解 forward-process、正负样本和 solver 解耦；
\item
  Ingredients / DexFlyWheel / RL-100：理解真实机器人系统和数据闭环；
\item
  OpenVLA / X-DiffVLA：理解 VLA、扩散动作头和跨 embodiment；
\item
  在仿真环境中比较 BC、Diffusion Policy、Flow-GRPO 风格更新和 NFT
  风格更新，再考虑真实机器人。
\end{enumerate}

\subsection{参考链接}\label{ux53c2ux8003ux94feux63a5}

\href{https://arxiv.org/abs/2505.05470}{Flow-GRPO 论文} ·
\href{https://github.com/yifan123/flow_grpo}{代码}；\href{https://arxiv.org/abs/2509.16117}{DiffusionNFT
论文} ·
\href{https://github.com/NVlabs/DiffusionNFT}{代码}；\href{https://arxiv.org/abs/2303.04137}{Diffusion
Policy}；\href{https://arxiv.org/abs/2510.14830}{RL-100}；\href{https://arxiv.org/abs/2509.23829}{DexFlyWheel}；\href{https://arxiv.org/abs/2406.09246}{OpenVLA}；\href{https://arxiv.org/abs/2605.25044}{X-DiffVLA}；\href{https://arxiv.org/abs/2311.01223}{Diffusion
Models for RL Survey}。

\end{document}

